\section{数值实验}\label{sec:numerics}

\begin{example}\label{ex:min-kernel}
考察 $L^2[0,1]$ 上的积分算子
\begin{equation}\label{eq:integral-operator}
  (Tf)(s)=\int_0^1\min\{s,t\}\,f(t)\,\mathrm{d}t .
\end{equation}
$T$ 紧且自伴非负. 解边值问题 $-\lambda u''=u$, $u(0)=u'(1)=0$, 得
\begin{equation}\label{eq:exact-eigen}
  \lambda_n=\frac{1}{(n-\frac12)^2\pi^2},\qquad e_n(s)=\sqrt2\,\sin\Bigl(\bigl(n-\tfrac12\bigr)\pi s\Bigr),\qquad n\in\N_+ .
\end{equation}
\end{example}

取中点求积公式, 节点 $s_i=(i-\frac12)h$, $h=1/N$, 则 \cref{eq:integral-operator} 化为矩阵特征值问题 $K_N v=\lambda v$, 其中 $(K_N)_{ij}=h\min\{s_i,s_j\}$. \Cref{alg:power} 给出求最大特征值的幂迭代, \cref{tab:eigen} 列出前五个特征值.

\begin{algorithm}[htbp]
\caption{带 Rayleigh 商的幂迭代}\label{alg:power}
\begin{algorithmic}[1]
\Require 对称矩阵 $K\in\R^{N\times N}$, 初值 $v_0$, 容差 $\varepsilon$
\State $v\gets v_0/\norm{v_0}$, $\mu\gets0$
\Repeat
  \State $w\gets Kv$
  \State $\mu_{\mathrm{old}}\gets\mu$, $\mu\gets v^{\top}w$
  \State $v\gets w/\norm{w}$
\Until{$\abs{\mu-\mu_{\mathrm{old}}}<\varepsilon\abs{\mu}$}
\State \Return $\mu$, $v$
\end{algorithmic}
\end{algorithm}

\begin{table}[htbp]
\centering
\caption{前五个特征值: Nyström 逼近与精确值}\label{tab:eigen}
\begin{tabular}{cccccc}
\toprule
$n$ & 精确值 $\lambda_n$ & $N=50$ & 相对误差 & $N=200$ & 相对误差 \\
\midrule
1 & 0.405285 & 0.405318 & $8.2\times10^{-5}$ & 0.405287 & $5.1\times10^{-6}$ \\
2 & 0.045032 & 0.045065 & $7.4\times10^{-4}$ & 0.045034 & $4.6\times10^{-5}$ \\
3 & 0.016211 & 0.016245 & $2.1\times10^{-3}$ & 0.016213 & $1.3\times10^{-4}$ \\
4 & 0.008271 & 0.008305 & $4.0\times10^{-3}$ & 0.008273 & $2.5\times10^{-4}$ \\
5 & 0.005004 & 0.005037 & $6.7\times10^{-3}$ & 0.005006 & $4.2\times10^{-4}$ \\
\bottomrule
\end{tabular}
\end{table}

网格加密 4 倍, 相对误差约降为原来的 $1/16$, 与中点公式的二阶精度一致. \Cref{fig:spectrum} 画出 $\sigma(T)$ 在实轴上的分布: 特征值以 $0$ 为唯一聚点, 与 \cref{thm:riesz-schauder} 相符.

\begin{figure}[htbp]
\centering
\begin{tikzpicture}[x=2.4cm,y=1cm,>=Stealth]
  % 横轴: log10(lambda), 取值 -2.6 .. 0
  \draw[->] (-3.1,0) -- (0.4,0) node[right] {$\log_{10}\lambda$};
  \foreach \t in {-3,-2,-1,0} \draw (\t,0.08) -- (\t,-0.08) node[below] {\small $\t$};
  \foreach \n/\v in {1/-0.3922,2/-1.3465,3/-1.7902,4/-2.0824,5/-2.3007,6/-2.4750,7/-2.6201,8/-2.7444,9/-2.8531,10/-2.9497}
    \fill[blue!60!black] (\v,0) circle (2pt);
  \node[above] at (-0.3922,0.15) {\small $\lambda_1$};
  \node[above] at (-1.3465,0.15) {\small $\lambda_2$};
  \node[above] at (-1.7902,0.15) {\small $\lambda_3$};
  \draw[<-,gray] (-3.05,0.25) -- (-3.05,0.9) node[above,align=center] {\small 聚点 $0$\\[-2pt]\small ($\log\lambda\to-\infty$)};
\end{tikzpicture}

\caption{积分算子 \cref{eq:integral-operator} 的谱 (对数坐标)}\label{fig:spectrum}
\end{figure}

计算 \cref{tab:eigen} 所用的代码如下.

\begin{lstlisting}[language=Python]
import numpy as np

def nystrom_eigs(N, k=5):
    h = 1.0 / N
    s = (np.arange(N) + 0.5) * h          # 中点节点
    K = np.minimum.outer(s, s) * h        # 离散核矩阵
    ev = np.linalg.eigvalsh(K)[::-1]      # 降序
    return ev[:k]

exact = 1 / (((np.arange(1, 6) - 0.5) * np.pi) ** 2)
for N in (50, 200):
    approx = nystrom_eigs(N)
    print(N, np.abs(approx - exact) / exact)
\end{lstlisting}
